\section{Types, Unification, and Constraints}
\label{sec:constraints}

There is one kind of object, a \emph{term} over the grammar's signature; one
semantic operation, \emph{ascription} (a check of a term against an expected
one); and one mechanism, \emph{unification}. Context lookup, type equality,
inclusion, and structural matching are all ascription discharged by unification.
The type constructors are the grammar's own productions, so there is no separate
algebra of types.

\subsection{Types as terms}

\begin{definition}[Type as term]
  \label{def:term-language}
  \code{def:term-language}
  The productions of the designated type fragment $\mathsf{Ty} \in N$ (written
  \texttt{(*)} on its production in the concrete grammar) form a
  \emph{signature} $\Sigma$: each production is a constructor whose arity is its
  right-hand nonterminals. A \emph{type} is a term in $T(\Sigma)$, a parse tree
  over $\mathsf{Ty}$. There is no privileged type algebra: $\to$, products, and
  unions are constructors only because the grammar writes them. Terminal leaves
  (type names) are recognized by regular terminals, so a type term is a tree
  with \emph{regular leaves}.
\end{definition}

\begin{remark}[The grammar is an order-sorted signature]
  \label{rem:sorts}
  The nonterminals are \emph{sorts}. An alternative whose body is a single
  nonterminal, $n \to n'$, is the subsort declaration $n' \leq n$, and a
  transparent production (Definition~\ref{def:normal-form}) is the injection it
  denotes, which is why collapsing it preserves meaning. The signature is
  therefore order-sorted~\cite{goguen1992order}: a hole at a position of sort
  $n$ ranges over terms of sort at most $n$
  (Definition~\ref{def:hole}), and a term \emph{is a type} exactly when its
  sort is below $\mathsf{Ty}$. This answers, inside the model, which terms may
  stand in type position; no construct outside the fragment can. Order-sorted
  unification stays unitary when sorts with a common subsort have a greatest
  one; this regularity condition is checkable at load time and belongs to the
  realizability gate.
\end{remark}

\subsection{Terms with holes}
\label{sec:patterns}

\begin{definition}[Hole]
  \label{def:hole}
  A \emph{hole} $?A, ?B, \dots$ is a variable standing for an unknown subterm. A
  hole at an internal position ranges over terms; a hole at a leaf position
  ranges over that leaf's regular language, so holes are sorted by the grammar's
  nonterminals (order-sorted, \cite{baader2001unification}). Two are
  distinguished: $\top$, a bare unconstrained hole, the most general type; and
  $\bot$, the absence of any term, a unification failure.
\end{definition}

\begin{definition}[Pattern]
  \label{def:pattern}
  \code{def:pattern}
  A \emph{pattern} is a term with holes, an element of $T(\Sigma, X)$. It is
  \emph{closed} when it has no holes (an ordinary type) and \emph{linear} when no
  hole name repeats. Structure lives in the tree; a repeated hole name is the
  only carrier of equality between positions.
\end{definition}

\subsection{Normal forms}

Equality is compared on normal forms, not on raw trees.

\begin{definition}[Normal form]
  \label{def:normal-form}
  A confluent, terminating rewrite $\mathrm{nf}$ sends each term to a unique
  normal form. The core assumes only the free part: collapse every
  \emph{transparent} node, a production that is delimiters plus a single
  semantic child. Redundant grouping then disappears, so $\mathrm{Int} \to
  (\mathrm{Int} \to \mathrm{Int})$ and $\mathrm{Int} \to \mathrm{Int} \to
  \mathrm{Int}$ are one term. Richer normalizers (declared equations
  \S\ref{sec:ambiguity}, or $\beta$-conversion for a dependent fragment) are
  named extensions; only confluence and termination are assumed here.
\end{definition}

\subsection{Unification}
\label{sec:unify}

\begin{definition}[Unification]
  \label{def:match}
  \code{def:match}
  $\mathrm{unify}(s, t)$ is the most general substitution $\theta$, if one
  exists, with $\theta\,\mathrm{nf}(s) = \theta\,\mathrm{nf}(t)$: structural
  equality of normal forms, leaf holes meeting their regular constraints. This is
  first-order syntactic unification~\cite{robinson1965,martelli1982}: $\top$ (a
  bare hole) unifies with everything, $\bot$ (a constructor clash, or an
  occurs-check failure) with nothing. \emph{Matching} is the one-sided case, a
  pattern against a closed type, and is the basis of ascription.
\end{definition}

\begin{lemma}[Unification is unitary and decidable]
  \label{lem:match-finite}
  \code{lem:match-finite}
  With the occurs-check on (finite terms), $\mathrm{unify}(s, t)$ either fails or
  returns a single most general unifier, unique up to renaming, decidably and in
  near-linear time. No branching arises.
\end{lemma}

\noindent The proof is standard~\cite{martelli1982} and deferred. Unitarity is
the point: the disambiguation that a string-matching model must perform does not
occur, because the parse already fixes which subterm each hole binds. Dropping
the occurs-check admits cyclic terms and gives rational-tree unification for
recursive types (\S\ref{sec:ambiguity}).

\subsection{Ascription: the only operation}

\begin{definition}[Ascription]
  \label{def:ascription}
  The single semantic primitive is the \emph{ascription}
  \[
    b : P,
  \]
  read ``the evidence term $\tau_b$ at $b$ unifies with the expected pattern
  $P$.'' Evaluating it emits $\mathrm{unify}(P, \tau_b)$
  (Definition~\ref{def:evidence}) and contributes the resulting substitution to
  the ambient solution. Read bidirectionally, $P$ is the checked type and
  $\tau_b$ the synthesized one.
\end{definition}

Every typing relation is an ascription, with shared holes the sole carrier of
equality:
\begin{itemize}[nosep]
  \item \emph{Equality} $\tau_1 = \tau_2$ is $\tau_1 : \tau_2$.
  \item \emph{Membership} $x \in \Gamma$ is $x : ?A$ solved against $\Gamma$.
  \item \emph{Inclusion} $\sigma \sqsubseteq \tau$ is subsumption: $\sigma$ is an
        instance of $\tau$ (the one-sided match). $\top$ is above all, $\bot$
        below all.
  \item A structural shape such as $\mathrm{Arrow}(?A, ?B)$ has no built-in
        meaning: $\mathrm{Arrow}$ is a constructor because the grammar wrote it,
        and $?A, ?B$ bind its subterms. Unifying $b : \mathrm{Arrow}(?A, ?B)$
        matches the constructor and binds each subterm by position. No
        decomposition step, and no notion of domain or codomain beyond which
        subterm each hole occupies.
\end{itemize}

\begin{remark}[Subtyping is the subsumption order]
  \label{rem:subtype-inclusion}
  A type at any stage is a term, ground or with holes. Refinement is term
  subsumption,
  \[
    \sigma \sqsubseteq \tau \iff \sigma \text{ is an instance of } \tau,
  \]
  so a hole-bearing (partial, unknown) type is \emph{more general} than its
  instances; $\top$ (a bare hole) is the top, $\bot$ (no unifier) the bottom.
  Unification is the meet: $\mathrm{unify}(\sigma, \tau)$ is their most general
  common instance. Regular sets survive at the leaves, where a leaf hole's value
  is the intersection of its regular constraints, so the prefix layer's regex
  machinery is the leaf case of one structural order. This is at once semantic
  subtyping (a more specific type is an instance) and the prefix order of
  incremental parsing (a partial type is more general, and input only
  instantiates it, Lemma~\ref{lem:hole-monotone}).

  The order is used in one direction. A more general (unknown) actual refutes
  but never confirms: a stable failure against a closed expectation forces
  $\mathsf{Lost}$ and stays failed under instantiation, while a success against
  a still-open actual is provisional. Provisionality costs nothing because
  verdicts are recomputed at every prefix: an edge satisfied today is
  re-checked under the further-instantiated actual tomorrow, and nothing
  irreversible leaves a node before it is exact (its right boundary fixed,
  \S\ref{sec:effects}). So the syntactic layer bounds the semantic one without
  fusing them.
\end{remark}

\subsection{Rules as constraint generators}

\begin{definition}[Rule]
  \label{def:rule}
  \code{def:rule}
  A rule attached to $n \in N$ is a pair $(\Phi, c)$ where the premises
  $\Phi = (b_1 : P_1, \dots, b_m : P_m)$ and the conclusion $c = (\,:\!P_c)$ are
  ascriptions over a shared hole alphabet. Firing the rule at an evidence node
  $\nu$ emits, into the constraint graph
  (Definition~\ref{def:constraint-graph}), one unification edge per premise plus
  the conclusion edge; holes shared between ascriptions become equality edges
  joining the corresponding subterms.
\end{definition}

A rule contributes no control flow of its own: it only \emph{adds edges}, and
shared holes are the joints. This is the constraint-generation reading of type
inference~\cite{PottierRemy2005}: a program is elaborated into a constraint, and
typing is solving it.

\subsection{Solving}

\begin{definition}[Solve]
  \label{def:solve}
  \code{def:solve}
  The solver closes the constraint graph under unification to a fixed point:
  select a pending edge, unify its endpoints, propagate the substitution to
  incident edges, and merge equal holes, until no edge can fire. Closure is
  monotone (substitutions only instantiate, equalities only merge), so the fixed
  point exists and is order-independent. The substrate is union-find over the
  term structure; constraint handling~\cite{Sneyers2009} drives the schedule, and
  the equational extension of \S\ref{sec:ambiguity} represents alternatives as an
  e-graph in the sense of equality saturation~\cite{Willsey2021}.
\end{definition}

\paragraph{Operational closure.}
The implementation realizes the closure bottom-up. A rule's holes are
$\alpha$-renamed at each node, so two firings never collide; the bindings a
node's solution fixes on \emph{foreign} variables (a child's, or one reaching
it through the context) are exported as the \emph{residual equations} of its
evidence, and a node first merges its children's equations into its own
substitution. A binding discovered deep in one subtree thus reaches every
other occurrence of the same metavariable at their least common ancestor:
closure of the constraint graph under unification, computed incrementally.
Metavariable identity is global. A context entry whose type is still a hole
stays one unknown across all its uses; only an explicit instantiation
($\mathsf{inst}$, the elimination of a polymorphic scheme's $\forall$)
freshens. Freshening an ordinary lookup instead would generalize the unknown
silently and accept programs with no type.

The three-valued verdict of \S\ref{sec:guarantees} is read off the closed graph:
\textsf{Satisfied} when every edge unifies and no frontier remains, \textsf{Lost}
when some edge fails \emph{stably}, and \textsf{Live} otherwise. A failure is
stable when no instantiation can repair it. Every failure of the free theory is
stable (a constructor clash or an occurs-violation survives substitution), but
modulo a rewrite theory a hole can block a redex, so normalization and
substitution no longer commute on open terms; there only a failure between
ground terms is stable, and a non-ground failure is postponed as a pending
constraint rather than pruned.

\subsection{Context and effects}
\label{sec:effects}

A rule may extend the context in two ways, and they are kept apart. A premise
\emph{setting} $\Gamma[b:P]$ is \emph{premise-local}: it extends only the
context under which that premise's own term is checked (the descent into the
child the premise binds), and it never reaches a sibling premise. The binder $b$
is itself such a child, resolved under the ambient context, so the setting does
not apply to the descent into $b$. A \emph{right-bound effect} $\nabla$, declared
by the conclusion, is the only thing that crosses to later siblings: it extends
the context they see; a declaration adding a binding is the canonical case. Only
an \emph{exact} node exports an effect, since a prefix node's right boundary is
not yet fixed, and effects compose left to right, applied by $\oplus$
(Definition~\ref{def:constraint-domain}). Operationally this needs two contexts
per node, not one: the context the node was \emph{predicted at} (its
resumption identity, fixed once) and the \emph{working} context that
accumulates left-sibling effects and feeds the next child. Keeping them distinct
is what lets an exported effect reach right siblings without moving where the
node itself closes.

\subsection{Monotonicity and realizability}
\label{sec:guarantees}

\begin{lemma}[Completion monotonicity]
  \label{lem:hole-monotone}
  \code{lem:hole-monotone}
  For an evidence node $e$, the set $\mathrm{comp}(e, s)$ of ground instances of
  its type term only shrinks under extension: $\mathrm{comp}(e, sr) \subseteq
  \mathrm{comp}(e, s)$.
\end{lemma}
\begin{proof}
  Reading input only instantiates holes and grows the context, never the
  reverse: substitution is monotone, so a bound hole stays bound, and a leaf
  hole's regular residual only shrinks under its
  derivative~\cite{owens2009regex}.
\end{proof}

By contraposition $\mathsf{Lost}$ is stable: once no instance satisfies a
constraint, none ever will.

\begin{theorem}[Solver correctness]
  \label{thm:solver-correct}
  \code{thm:solver-correct}
  For the syntactic core the solver verdict on $\mathcal{G}(s)$ equals
  $\mathit{eval}(\mathcal{G}(s))$ (Definition~\ref{def:eval}) at every reachable
  $s$.
\end{theorem}

\noindent\emph{Proof (sketch; full proof deferred).} Induct on the AST. Each
non-$\mathsf{Lost}$ ascription has a satisfying continuation by productivity, and
$\mathsf{Lost}$-stability (Lemma~\ref{lem:hole-monotone}) forbids recovery;
unitarity (Lemma~\ref{lem:match-finite}) makes the per-node unifier unique, so the
join over children matches the open/closed split of
Definition~\ref{def:open-closed}. The equational extension of
\S\ref{sec:ambiguity} is the deferred case.

\subsection{Where multiplicity lives: beyond syntactic unification}
\label{sec:ambiguity}

The syntactic core is unitary, so rule firing never branches and there is no
disambiguation problem: the parse fixes the structure each hole binds.
Multiplicity reappears only when unification runs \emph{modulo} a theory the
grammar declares, each an opt-in extension rather than engine machinery:
\begin{enumerate}[nosep]
  \item \textbf{Equational equivalences.} Declaring $A \mid B \equiv B \mid A$
        and the like makes equality unification modulo an equational
        theory~\cite{baader2001unification}, which is finitary but not unitary
        (associative-commutative unification is the classic
        case~\cite{stickel1981}). The e-graph of Definition~\ref{def:solve}
        represents the alternatives without committing.
  \item \textbf{Recursive types.} Dropping the occurs-check admits cyclic terms;
        unification becomes rational-tree
        unification~\cite{colmerauer1982,Tiurin2024}, still decidable.
  \item \textbf{Conversion.} For a dependent fragment the normal form of
        Definition~\ref{def:normal-form} is $\beta$-normal form and equality is
        conversion, lifting unification to higher-order pattern
        unification~\cite{miller1991}. Decidability rests on strong
        normalization. This is the regime of a small dependent
        kernel~\cite{harper1993lf} and the eventual target of the framework.
\end{enumerate}
The core stays first-order and unitary; the cost is exactly the equational power
a grammar asks for.
